\documentclass[11pt]{article}
\usepackage[T1]{fontenc}
\usepackage{lmodern,amsmath,amssymb,booktabs,xurl}
\usepackage[margin=0.8in]{geometry}
\usepackage[colorlinks=true,urlcolor=blue]{hyperref}
\setlength{\parindent}{0pt}
\setlength{\parskip}{3pt}
\setlength{\emergencystretch}{2em}
\newcommand{\heading}[1]{\vspace{4pt}\textbf{#1}\par\nopagebreak}
\title{Local Models and Joint Queries\\\large What coherent revision requires}
\author{}
\date{}
\begin{document}
\begin{center}{\LARGE Local Models and Joint Queries}\\[3pt]{\large What coherent revision requires}\end{center}
\textbf{A reasoner can work with local probability models and have no largest
working context. For each joint question it answers, it must justify the
relevant joint interpretation---or report what remains undetermined.}

\heading{1. Three meanings of ``global''}
A \emph{working context} contains a set of quantities and a probability model
for them. A \emph{joint model} connects quantities that a question uses
together. A \emph{global extension} is one probability measure whose marginals
match every local model under consideration.

A joint model can be global relative to one query and still be small relative
to the whole system. Its existence does not require a master table, a largest
context, or exhaustive inference over everything the system knows. An abstract
global distribution, however, is still a global distribution mathematically;
``abstract'' does not mean it has disappeared.

\heading{2. No largest context is required}
Take all finite subsets of an infinite variable set. Any two have a common
enlargement, their union. Every context can grow, so none is greatest.
Independent fair coins give a coherent example: each finite context has its
finite product distribution, and there is also an infinite product measure.
Thus \emph{no largest finite context} and \emph{a global extension exists} can
both be true.

More generally, fix finite, nonempty value sets for the variables. If a joint
law $\mu_C$ is supplied for \emph{every finite variable set} $C$, and
\[
  (\text{restriction from }D\text{ to }C)_*\mu_D=\mu_C
  \qquad(C\subseteq D),
\]
then there is a unique global probability measure on the product
sigma-algebra with these finite marginals. This is the Kolmogorov extension
theorem in its finite-state form. It establishes existence, not an efficient
global algorithm.

The words \emph{every finite variable set} matter. Agreement between the
overlaps of a few local models is a weaker condition. For a partial collection of prescribed
finite marginals, if \emph{every finite collection of constraints}
has a common realization, compactness gives at least one global extension;
it need not be unique. Checking only pairs of constraints does not suffice.
These statements assume a fixed, commonly interpreted variable family.

\heading{3. A local joint is enough, but its dependence must be specified}
Suppose $X$ and $Y$ are both fair bits. Their separate laws allow:
\begin{center}
\begin{tabular}{lcc}
\toprule
Joint relationship & $P(X=1,Y=1)$ & $P(X=1\mid Y=1)$\\
\midrule
Always equal & $1/2$ & $1$\\
Independent & $1/4$ & $1/2$\\
Always opposite & $0$ & $0$\\
\bottomrule
\end{tabular}
\end{center}
A local joint model for $(X,Y)$ answers the revision question. The marginals
alone do not select it. Taking their product is a legitimate modeling choice
that asserts independence. Distinct variable names or evidence IDs do not
establish that assumption.

\heading{4. What failure of a global extension actually shows}
Consider three fair-bit pair models for $(X,Y)$, $(Y,Z)$ and $(X,Z)$. In each,
assign probability $1/16$ to each equal outcome and $7/16$ to each unequal
outcome. All single-bit overlaps agree. Yet every Boolean triple has at least
one equal pair, so any triple law must satisfy
\[
 P(X=Y)+P(Y=Z)+P(X=Z)\geq1.
\]
The proposed pair laws give $3/8$. No triple joint can realize all three.

This disproves the claim that overlap agreement always permits joining local
models. It does not disprove the extension theorem: the required law on
$\{X,Y,Z\}$ is precisely what is missing. It also does not make the local
calculations useless. They remain valid within their declared contexts;
combining all three as one classical joint interpretation is unavailable.
Treating $X$ in one chart and $X$ in another as different variables is a
legitimate move, but it changes the problem: the identification must be
dropped explicitly, not silently.

For \emph{two} finite models $p(a,s)$ and $q(s,b)$ with the same overlap
marginal $r(s)$ and no extra cross-constraints, one possible joining is
\[
 \mu(a,s,b)=\begin{cases}p(a,s)q(s,b)/r(s),&r(s)>0,\\0,&r(s)=0.\end{cases}
\]
This conditional-independence bridge preserves both supplied marginals.
For the triangle, joining the first two charts gives
$P(X=1\mid Z=1)=25/32$, whereas the third gives $1/8$.
A path label records which model was used; it does not reconcile the answers.

\heading{5. Say what the information supports}
Let $\mathcal K$ be the set of joint models consistent with the declared
constraints on a requested context.
\begin{itemize}
\setlength{\itemsep}{0pt}
\item If $\mathcal K$ is nonempty and all its models agree on the query, return that answer,
      even when the full joint model is not unique.
\item If answers vary, return bounds or expose the assumption selecting one
      model. For two fair bits, the conjunction lies in $[0,1/2]$ and the
      conditional in the table lies in $[0,1]$.
\item If $\mathcal K$ is empty, report incompatible joint constraints.
      Preserve useful context-specific results, or explicitly repair/select
      the assumptions. Do not manufacture a probability for the empty model set.
\end{itemize}
For conditioning on $E$, evaluate $P(A\mid E)$ only for admissible models
with $P(E)>0$. Impossible evidence, unknown dependence, and incompatible
constraints are different outcomes. The failure of one joint request need not
silence the rest of the reasoner.

A prior over candidate models is one way to supply joint models. When the
candidates describe the same observables, their mixture is a joint model and
evidence revises the weights, $\pi(m\mid e)\propto\pi(m)\,P(e\mid m)$:
likelihoods $9/10$ and $1/10$ turn equal weights into $9/10$ and $1/10$. No
largest context is needed, only an agreed meaning for the shared variables.

\heading{6. The operational contract and checked result}
Retain the model, evidence identities, query interpretation, and the
justification for each coupling or likelihood product. Work in any context sufficient for
the current question; enlarge it when needed. Different larger contexts give
the same revised answer if they agree on the joint query--evidence law and
use the same likelihood. New evidence may change the answer.

\textbf{Checked in Lean:} the finite query--evidence reduction and normalizer
agreement, directed finite contexts without a greatest member, the two-coupling
revision control, the strictly positive triangle obstruction, the two-chart
bridge ($25/32$ against $1/8$), and the model-prior update. See the Mettapedia module
\texttt{ProbabilityTheory/\allowbreak BayesianInference/\allowbreak LocalContexts.lean}.
The general extension/compactness statements above are classical mathematical
background, not new Lean results in that module.
\end{document}
